\section{Upper bounds of MED-C, Simulation, and Implications}~\label{sec:maed_and_exp}

This section discusses the minimal embeddable dimension in the centroid setting (MED-C).
We will see that MED-C grows logarithmically with the universe's size $m$ in both theory and experiments.

\subsection{Upper bounds}

For inner product, cosine similarity, and Euclidean distance, we use the probabilistic method as a common proving strategy to prove the logarithmic upper bound. In short, we sample $m$ element embeddings from a probabilistic distribution in $\real^n$ and find the conditions on $k$, $m$, and $n$ such that the desired $k$-centroid shattering configuration occurs with positive probability. The random vectors here are part of the probabilistic construction and do not directly link to the simulation.

\begin{theorem}\label{thm:medc-linear}
    $\medc(m, k; s_{\rm linear})=O(k^2\log m)$
\end{theorem}
\noindent\textit{Proof sketch:} We consider the probabilistic method, where $v_1, \dots, v_m \sim \mathcal{N}(0, 1/n)$ in $\real^n$, we show there is positive probability such that for any subset $S$, $|S|\leq k$, $v_1\in S$ and $v_2\notin S$
$ \langle v_1, \sum_{u\in S} u \rangle > \langle v_2 ,  \sum_{u\in S} u \rangle$ under the condition of $n>Ck^2\log m$. The full proof can be found in Appendix~\ref{app:proof:thm:maed-linear}.

Following a similar probabilistic construction, we can also prove that $\fcos$ and $\fl{2}$ share the same upper bound. The proofs are also in the Appendix~\ref{app:proof:thm:maed-linear}.
\begin{theorem}\label{thm:medc-cos}
$\medc(m, k; \fcos)=O(k^2\log m).$
\end{theorem}
\begin{theorem}\label{thm:medc-l2}
$\medc(m, k; \fl{2})=O(k^2\log m).$
\end{theorem}

Although the $O(k^2\log m)$ bound may not be tight, it still suffices to make significant implications. An easy lower bound for MED-C is the lower bound of MED, which is $\Omega(k)$. We can see that the gap between the lower and upper bounds is a factor of $O(\log m)\times O(k)$. We discuss those two factors separately. For the factor of $O(\log m)$, it still demonstrates that a sublinear growth for MED-C with the size $m$ of the universe is feasible. For more complex neural set embedding settings, such as MED-N discussed in Section~\ref{sec:maed}, a similar upper bound also applies. For the additional factor $O(k)$, we would like to convince the readers that $O(k\log m)$ is not very easy to achieve, so the $O(k^2 \log m)$ is not very bad. To the best of the authors' knowledge, the most similar (but not exact) problem setting to MED-C is 1-bit compress sensing~\citep{aksoylar2014information}, with a lower bound of $\Omega(k\log \frac{m}{k}) $, which cancels the additional $O(k)$ factor (compare $O(k^2\log m)$ for MED-C with $\Theta(k)$ for MED) but requires an additional sparse recovery algorithm rather than simply comparing scores independently. This related work suggests that the extra cost should be paid to achieve the bound $O(k\log m)$ of MED-C.


\subsection{Numerical simulation}

The situation where $k \ll m$ is of particular interest because there are many documents to be retrieved, and only a very small subset is required in a real-world scenario. Therefore, numerical simulation typically focuses on large $m$ and small $k$ to demonstrate the effectiveness of the bounds.
In the numerical simulation, we use "critical" to denote the largest possible number of objects for a given dimension, or the smallest dimension required to accommodate a specified number of objects.

\noindent\textbf{Experiment settings.} We use optimize $m$ embeddings randomly initialized in by standard normal distribution in $\real^d$. The objective function is a hinge loss computed over all pairwise positive and negative pairs across the top-$k$ queries. We use the Adam optimizer~\citep{kingma2014adam} with a learning rate of 1. The optimization stops either after 1,000 steps or after every query finds its perfect answers. To facilitate full reproducibility, we list our Python code in the Appendix to disclose other details.

\noindent\textbf{Comparison with baselines}
Previous work~\citep{weller2025theoretical} didn't provide a quantitative relation in their theory, but instead offered an empirical relation via optimization of free embeddings, which corresponds to the standard setting in this paper. They search for the critical cardinality of the universe $m^*$ in a fixed dimension $d$ for the $k=2$ case, and established the following empirical results. We list their curve in Equation~\ref{eq:baseline}.
\begin{align}\label{eq:baseline}
    m_{\rm WBNL}(d) = -10.5322 + 4.0309 d + 0.0520d^2 + 0.0037 d ^3.
\end{align}
This result yields a pessimistic implication: the largest number of objects we can support grows with $d$ only cubically. However, our  bounds in MED-C suggests that the number of objects grows exponentially with $d$.

\begin{figure*}[t]
    \centering
    \begin{subfigure}[t]{0.48\textwidth}
        \centering
        \includegraphics[width=\textwidth]{figure/compare_plot1.pdf}  % or your first page/image
        \caption{}  % empty caption here (required technically, but suppressed)
        \label{fig:first}
    \end{subfigure}
    \hfill  % maximizes separation between the two
    \begin{subfigure}[t]{0.48\textwidth}
        \centering
        \includegraphics[width=\textwidth]{figure/compare_plot2.pdf}  % or your second page/image
        \caption{}  % empty caption here
        \label{fig:second}
    \end{subfigure}

    % Single shared caption for both subfigures
    \caption{(a) The comparison of the growth of the observed supported number of points from cyclic-polytope and centroid GD upper-bound witnesses against the curve fitted in Equation~\ref{eq:baseline}.; (b) The comparison of the witness dimensions from the same grid scan against the curve fitted in Equation~\ref{eq:baseline}, with the $x$ axis plotted in a log scale.}
    \label{fig:two-pages-together}
\end{figure*}

It is a crucial question: how can the critical points identified by simulation in the centroid setting be compared with those found in the free embedding optimization setting~\citep{weller2025theoretical}.
The goal of \textbf{numerical simulation} is to estimate the minimum dimension such that $m$ elements can be $k$-shattered in some sense (or the maximum number of elements can be located in a given dimension). Therefore, all critical points found by any numerical simulations only suggest an upper bound of the minimal dimension (or lower bound of the maximum number of elements). The goal of \textbf{this paper} is to show that the MED (or the lower bound) grows slowly enough so that the fundamental property of the geometric space does not bottleneck the embedding-based retrieval. Therefore, results from simulations, whether computed from a centroid setting or a free embedding optimization setting, can serve as empirical evidence of the MED upper bound.

To compare with their empirical results, we ran the experiments with $k=2$ using both the cyclic-polytope construction and the centroid setting. We plot the results in Figure~\ref{fig:first} and Figure~\ref{fig:second}. Figure~\ref{fig:first} suggests that the number of supported points found by these upper-bound witnesses surpasses the curve fitted by~\citet{weller2025theoretical} easily. From Fig~\ref{fig:second}, we can see from the results that our theory of MED-C agrees well with the empirical log-linear fitting, which is of a lower complexity family than the fitted result in~\citet{weller2025theoretical}.



\subsection{Practical implication}

\noindent\textbf{Regarding the bottleneck of the embedding-based retrieval.} The empirical results above suggested that the actual limitation of embedding-based retrieval, i.e., the number of objects that can be embedded, actually grows at least exponentially with the dimension. This contradicts the previous claims made by~\citet{weller2025theoretical}: The stagnation of the performance in the embedding-based system does NOT necessarily come from the fundamental geometric property of the vector space. The empirical findings resonated with the conclusion from $\Theta(k)$ bounds in a weaker form: we can observe the exponential growth of critical points and logarithmic growth of MED, but not a vertical line or horizontal line.

\noindent\textbf{Why numerical simulation in centroid setting achieves stronger points than free embedding optimization?}
It is also a little bit counterintuitive that, optimizing with $\Theta(\binom{m}{k})$, more degrees of freedom actually leads to weaker results. Specifically, one should expect the optimums of the free embedding optimization and centroid setting to follow the following inequality:
\begin{align*}
    &\overbrace{\min_{\vw_S: S\in \collection_k}}^{\text{additional variables}} \min_{\vx: x\in X}  L_{\rm Free Embedding} \leq \min_{\vx: x\in X} L_{\rm Centroid},
\end{align*}
where we define hinge losses $L_{\rm Free Embedding} := \sum_{S \in \collection_k,x\in S,y\notin S}\max(0, s(\vy, \vw_S) - s(\vx, \vw_S))$ and $L_{\rm Centroid} :=  \sum_{S \in \collection_k,x\in S,y\notin S}\max(0, s(\vy, \frac{1}{|S|}\sum_{x\in S} \vx) - s(\vx, \frac{1}{|S|}\sum_{x\in S} \vx))$ are optimization objectives in free-embedding setting and centroid setting, respectively.

However, the simulation empirically reveals the inequality in the other direction. It is another perfect example of how optimization affects the final results and further distorts the empirical findings.
